\documentclass[12pt]{article}
\usepackage[margin=1in]{geometry}
\usepackage{setspace}

\setstretch{1.5}

\begin{document}

\begin{center}
\textbf{Individual Contribution Report}
\end{center}

\textbf{Name:} Tuğçe Yücel \\
\textbf{Student ID:} 2022400105 \\
\textbf{Group Number:} 31 \\

\section*{Tasks \& Contributions}

In this project, I was responsible for designing several core components of the database system both at the conceptual (ER design) and logical (SQL implementation) levels.

During the conceptual design phase (Part 1), I designed the \textit{Person}, \textit{Club}, \textit{Contract}, and \textit{Transfer Record} entities in the ER diagram. This involved identifying their attributes, defining primary keys, determining appropriate relationships, and applying ER constraints such as key constraints and participation constraints. Special attention was given to modeling \textit{Contract} and \textit{Transfer Record}, since they represent historical records that must be permanently stored in the system.

For the logical design phase (Part 2), I implemented the SQL DDL statements corresponding to these components of the database. I defined relational tables, primary keys, foreign keys, and additional integrity constraints such as \texttt{CHECK}, \texttt{NOT NULL}, and \texttt{UNIQUE}. Furthermore, I prepared the Part 2 report by documenting the relational schema and explaining the mapping from the ER design to the SQL implementation.

\section*{Collaboration \& Challenges}

Our team worked collaboratively throughout the design process. Initially, each team member focused on designing different parts of the ER model individually. After completing our respective sections, we combined them into a single diagram and reviewed the overall structure together.

One of the main challenges was determining the most appropriate way to represent certain concepts in the ER diagram. For example, we needed to carefully decide whether certain structures should be modeled as entities or relationships, and how to represent complex constraints related to contracts and player transfers. Through group discussions and multiple revisions of the diagram, we refined the design and resolved these issues.

\section*{Self-Assessment}

This project significantly improved my understanding of database design and implementation. While working on the ER diagram, I gained valuable experience in making design decisions and understanding how different entities and relationships should be structured within a database model.

In addition, writing the SQL DDL statements helped me develop stronger practical skills in SQL. By implementing constraints and translating the conceptual model into relational tables, I gained a deeper understanding of how theoretical database design principles are applied in practice. Overall, the project strengthened both my conceptual understanding of database systems and my technical skills in SQL.

\section*{Use of AI Tools}

AI tools were used occasionally as supportive assistants during the project. They were mainly consulted when we needed guidance in evaluating different modeling approaches or when we wanted to verify the correctness of certain database design decisions.

These tools helped us explore alternative solutions and better understand the implications of different design choices. However, the final ER diagram, SQL implementation, and documentation were produced by the team based on our own understanding of the project requirements.

\end{document}